\section{Introducción general}
\subsection{Sistema de detección de apneas y variaciones de oxigenación durante el Sueño}
El sistema portátil de monitoreo respiratorio y de saturación de oxígeno desarrollado en este trabajo cuenta con:
\begin{itemize}
    \item Una placa con microcontrolador \textit{STM32 F103RB}, encargada de la adquisición, procesamiento digital de señales y control general del sistema.
    \item Un sensor óptico \textit{MAX30102}, que permite medir señales PPG mediante leds rojo e infrarrojo para el cálculo de $SpO_{2}$ y la estimación de la frecuencia respiratoria.
    \item Un módulo de comunicación Bluetooth \textit{HM-10}, que posibilita la transmisión inalámbrica de los datos hacia un dispositivo móvil para su visualización en tiempo real.
    \item Un display, que muestra datos por pantalla en tiempo real.
    \item Un sistema de alarmas locales compuesto por un buzzer y leds indicadores, que alertan al usuario ante condiciones críticas de saturación de oxígeno.
    \item Procesamiento digital de la señal implementado en lenguaje C.
    \item Lógica de control que activa alarmas cuando la saturación de oxígeno cae por debajo de los umbrales configurados.
    \item Modos de funcionamiento configurables (modo adulto, modo niño).
\begin{figure} [H]
    \centering
    \includegraphics[width=0.6 \linewidth]{diagrama_de_bloques.png}
    \caption{Diagrama de bloques del sistema}
    \label{fig:diag_bloques}
\end{figure}
\end{itemize}

\subsection{Análisis de sistemas similares al desarrollado}
La inspiración para el desarrollo de este proyecto se basa en el funcionamiento de un oxímetro de pulso, el cual utiliza señales PPG obtenidas mediante la emisión de luz roja e infrarroja para estimar la saturación de oxígeno en sangre.
Tomando como referencia este principio de funcionamiento, el presente proyecto propone extender sus capacidades, con el objetivo de desarrollar un sistema de monitoreo completo durante el sueño. Además de estimar la saturación de oxígeno,permite obtener y supervisar otros parámetros fisiológicos relevantes, como la frecuencia cardíaca, la frecuencia respiratoria y la detección de episodios de apnea, proporcionando así una herramienta más completa para el seguimiento del estado fisiológico del usuario durante el descanso.
\begin{figure}
    \centering
    \includegraphics[width=0.4\linewidth]{oximetro.png}
    \caption{Oxímetro de pulso}
    \label{fig:oximetro}
\end{figure}
